\section{Особенности реализации}

\subsection{Среда реализации}

Экспертная система реализована в среде CLIPS. Эта среда подходит для выбранной продукционной модели, так как позволяет описывать знания через правила, хранить ответы пользователя в рабочей памяти и запускать прямой логический вывод. Математическая модель дерева решений описана в разделе~\ref{sec:decision-tree-model}, а в данном разделе рассматривается её программная реализация.

\subsection{Шаблоны}

Рабочая память системы хранит ответы пользователя в однотипных фактах \texttt{known}. Структура такого факта приведена в листинге~\ref{lst:known-template}: поле \texttt{attribute} задаёт имя проверяемого признака, а поле \texttt{value} хранит бинарное значение \texttt{yes} или \texttt{no}.

\begin{listing}[H]
\caption{Шаблон факта \texttt{known}}
\label{lst:known-template}
\inputminted[
    linenos,
    numbersep=10pt,
    firstline=1,
    lastline=4,
    fontsize=\small
]{text}{../fonts_advisor.clp}
\end{listing}

\subsection{Функции}

Для обработки ввода используется функция \texttt{ask-choice-1-2}, приведённая в листинге~\ref{lst:ask-choice}. Она повторяет запрос до тех пор, пока пользователь не введёт значение \texttt{1} или \texttt{2}. Такой вариант выбран вместо свободного ввода строк, чтобы избежать неоднозначных ответов и опечаток.

\begin{listing}[H]
\caption{Функция проверки бинарного выбора}
\label{lst:ask-choice}
\inputminted[
    linenos,
    numbersep=10pt,
    firstline=7,
    lastline=20,
    fontsize=\small
]{text}{../fonts_advisor.clp}
\end{listing}

Универсальная функция вопроса \texttt{ask-binary} показана в листинге~\ref{lst:ask-binary}. Она выводит текст вопроса, вызывает функцию проверки выбора из листинга~\ref{lst:ask-choice}, после чего добавляет в рабочую память факт \texttt{known} с соответствующим значением.

\begin{listing}[H]
\caption{Функция задания бинарного вопроса}
\label{lst:ask-binary}
\inputminted[
    linenos,
    numbersep=10pt,
    firstline=24,
    lastline=38,
    fontsize=\small
]{text}{../fonts_advisor.clp}
\end{listing}

Функция \texttt{print-result}, представленная в листинге~\ref{lst:print-result}, используется всеми листовыми правилами. Она печатает номер листа дерева и список рекомендаций, после чего завершает работу системы командой \texttt{halt}.

\begin{listing}[H]
\caption{Функция вывода результата}
\label{lst:print-result}
\inputminted[
    linenos,
    numbersep=10pt,
    firstline=41,
    lastline=50,
    fontsize=\small
]{text}{../fonts_advisor.clp}
\end{listing}

\subsection{Факты}

Факты создаются только после ответа пользователя на вопрос. Например, если пользователь подтверждает наличие кириллицы, функция из листинга~\ref{lst:ask-binary} добавляет факт \texttt{(known (attribute cyrillic) (value yes))}. Если ответ отрицательный, создаётся аналогичный факт со значением \texttt{no}.

Такая организация делает правила независимыми от порядка ввода: каждое правило проверяет только те факты, которые нужны для соответствующей ветви дерева. Пример правила-вопроса приведён в листинге~\ref{lst:ask-cyrillic}. Оно срабатывает только после подтверждения наличия дополнительного текста и задаёт следующий вопрос о кириллице.

\begin{listing}[H]
\caption{Правило-вопрос о наличии кириллицы}
\label{lst:ask-cyrillic}
\inputminted[
    linenos,
    numbersep=10pt,
    firstline=62,
    lastline=68,
    fontsize=\small
]{text}{../fonts_advisor.clp}
\end{listing}

Листовые правила проверяют полное сочетание признаков и вызывают функцию вывода результата из листинга~\ref{lst:print-result}. Пример такого правила приведён в листинге~\ref{lst:node-16}; оно соответствует одному из конечных листьев дерева на рисунке~\ref{fig:decision-tree}.

\begin{listing}[H]
\caption{Пример листового правила}
\label{lst:node-16}
\inputminted[
    linenos,
    numbersep=10pt,
    firstline=361,
    lastline=376,
    fontsize=\small
]{text}{../fonts_advisor.clp}
\end{listing}
